\documentclass[10pt]{article}

\usepackage[latin1]{inputenc}
\usepackage{amsmath, amssymb, amsfonts, amsthm}
\usepackage{upgreek}
\usepackage{amsthm}
\usepackage{fullpage}
\usepackage{graphicx}
\usepackage{cancel}
\usepackage{subfigure}
\usepackage{mathrsfs}
%\usepackage{outlines}
\usepackage[font={sf,it}, labelfont={sf,bf}, labelsep=space, belowskip=5pt]{caption}
\usepackage{hyperref}
%\usepackage{minted}
\usepackage{xcolor}

\usepackage{fancyhdr}
\usepackage[title]{appendix}

\DeclareMathOperator{\sgn}{sgn}
\DeclareMathOperator{\tr}{tr}

\pagestyle{fancy}
\headheight 24pt
\headsep    12pt
\lhead{Tensor Flattening}
\fancyfoot[C]{} % hide the default page number at the bottom
\lfoot{}
\rfoot{\thepage}
\renewcommand{\headrulewidth}{0.4pt}
\renewcommand\footrulewidth{0.4pt}
\providecommand{\abs}[1]{\lvert#1\rvert}
\providecommand{\norm}[1]{\lVert#1\rVert}
\providecommand{\dx}{\, \mathrm{d}x}
% \providecommand{\vint}[2]{\int_{#1} \! #2 \, \mathrm{d}x}
% \providecommand{\sint}[2]{\int_{\partial #1} \! #2 \, \mathrm{d}A}
\renewcommand{\div}{\nabla \cdot}
\providecommand{\shape}{\Omega(p)}
\providecommand{\boundary}{\partial \shape}
\providecommand{\vint}[1]{\int_{\shape} \! #1 \, \mathrm{d}x}
\providecommand{\sint}[1]{\int_{\boundary} \! #1 \, \mathrm{d}A}
\providecommand{\pder}[2]{\frac{\partial #1}{\partial #2}}
\providecommand{\tder}[2]{\frac{\mathrm{d} #1}{\mathrm{d} #2}}
\providecommand{\evalat}[2]{\left.#1\right|_{#2}}
\providecommand{\e}{\epsilon}
\newcommand{\defeq}{\vcentcolon=}
\newtheorem{lemma}{Lemma}

\makeatletter
\usepackage{mathtools}
\newcases{mycases}{\quad}{%
  \hfil$\m@th\displaystyle{##}$}{$\m@th\displaystyle{##}$\hfil}{\lbrace}{.}
\makeatother

\title{Tensor Flattening}
\author{Julian Panetta}

% BEGIN DOCUMENT
\begin{document}

\maketitle

We discuss how symmetric rank 2 and rank 4 tensors are flattened into vectors
and matrices so that contractions can be implemented efficiently as dot products
and matrix multiplies. We use summation notation for our tensors, where
summation over repeated indices is implied. The range of indices is $0..2$ in
3D and $0..1$ in 2D.

\section{Voigt Notation}
Our approach is to use Voigt notation with zero indexing.
\subsection{Symmetric Rank 2 Tensors}
In this notation, symmetric rank 2 tensors are flattened into vectors in the order:
$$
\text{3D:} \quad
\begin{pmatrix}
    0 & 5 & 4 \\
    \textcolor{lightgray}{5} & 1 & 3 \\
    \textcolor{lightgray}{4} & \textcolor{lightgray}{3} & 2
\end{pmatrix} \quad
\text{2D:} \quad
\begin{pmatrix}
    0 & 3 \\
    \textcolor{lightgray}{3} & 1
\end{pmatrix}.
$$
For example, strain tensors are flattened as:
$$
\text{3D:} \quad
\begin{pmatrix}
    \e_{xx} & \e_{xy} & \e_{xz} \\
    \e_{xy} & \e_{yy} & \e_{yz} \\
    \e_{xz} & \e_{yz} & \e_{zz}
\end{pmatrix} \mapsto
\begin{pmatrix}
    \e_{xx} \\
    \e_{yy} \\
    \e_{zz} \\
    \e_{yz} \\
    \e_{xz} \\
    \e_{xy}
\end{pmatrix}, \quad \quad
\text{2D:} \quad
\begin{pmatrix}
    \e_{xx} & \e_{xy}\\
    \e_{xy} & \e_{yy}\\
\end{pmatrix} \mapsto
\begin{pmatrix}
    \e_{xx} \\
    \e_{yy} \\
    \e_{xy}
\end{pmatrix} \quad
$$
This flattening is implemented by the following formula for flattened index
$\overline{ij}$ in $N$ dimensions:
$$
\overline{ij}(i, j) = \begin{cases}
    i & \text{if } i = j \\
    \frac{N (N + 1)}{2} - \left(i + \frac{j (j - 1)}{2} + 1\right) & i < j \\
    \overline{ij}(j, i)            & i > j
\end{cases} .
$$
We can also go the other direction by solving for the upper triangle $(i, j)$ associated with $\overline{ij}$.
This is done by picking the largest $j$ possible (to ensure we're in the upper
triangle), and then solving for the $i$ index:
$$
j(\overline{ij}) = \begin{cases}
    \overline{ij}  & \text{if } \overline{ij} < N \\
    \left\lfloor \frac{1 + \sqrt{1 + 8 \left(\frac{N (N + 1)}{2} - 1 - \overline{ij} \right)}}{2} \right\rfloor & \text{otherwise}
\end{cases} \quad \quad
i(\overline{ij}) = \begin{cases}
    \overline{ij}  & \text{if } \overline{ij} < N \\
    \frac{N (N + 1)}{2} - 1 - \overline{ij} - \frac{j(\overline{ij}) (j(\overline{ij}) - 1)}{2} & \text{otherwise}
\end{cases}.
$$
\subsection{Symmetric Rank 4 Tensors}
The symmetric rank 2 tensor flattening induces a flattening of symmetric rank 4
tensors that act on symmetric rank 2 tensors by double contraction. We consider
tensors with the symmetries of an elasticity tensor (i.e. major and minor
symmetries):
\begin{equation}
\label{eqn:symmetry}
E_{ijkl} = E_{jikl} = E_{ijlk} = E_{klij}.
\end{equation}
The idea is to flatten this tensor into a matrix so that matrix multiplication
can implement double contraction on both the left and the right: $E : \e =
E_{ijkl} \e_{kl}$ and $\e : E = \e_{ij} E_{ijkl}$. Since $\e_{kl}$ is flattened
into vector entry $\overline{kl}$, we want entries $E_{**kl}$ to end up in
matrix column $\overline{kl}$ so that entries are paired properly for
the left double contraction. Likewise, the right double contraction requires
$E_{ij**}$ to end up in matrix row $\overline{ij}$. Accordingly, we flatten
rank 4 tensors $E_{ijkl}$ into the matrix:
$$ D_{\overline{ij}\,\overline{kl}} = E_{ijkl} \quad \quad \text{where } i = i(\overline{ij}),\ j = j(\overline{ij}),\ k = i(\overline{kl}),\ l = j(\overline{kl})$$
The minor symmetries ($i\leftrightarrow j$ and $k \leftrightarrow l$) of tensor
$E_{ijkl}$ mean no information is lost in this conversion, and
the major symmetry $ij \leftrightarrow kl$ means $D$ is
a symmetric matrix.

\subsection{Examples}
A general rank 4 tensor looks like:
$$
\text{3D:}\quad
E_{ijkl} =
\left(
\begin{array}{ccc}
 \left(
\begin{array}{ccc}
 E_{0000} & E_{0001} & E_{0002} \\
 E_{0010} & E_{0011} & E_{0012} \\
 E_{0020} & E_{0021} & E_{0022} \\
\end{array}
\right) & \left(
\begin{array}{ccc}
 E_{0100} & E_{0101} & E_{0102} \\
 E_{0110} & E_{0111} & E_{0112} \\
 E_{0120} & E_{0121} & E_{0122} \\
\end{array}
\right) & \left(
\begin{array}{ccc}
 E_{0200} & E_{0201} & E_{0202} \\
 E_{0210} & E_{0211} & E_{0212} \\
 E_{0220} & E_{0221} & E_{0222} \\
\end{array}
\right) \\
 \left(
\begin{array}{ccc}
 E_{1000} & E_{1001} & E_{1002} \\
 E_{1010} & E_{1011} & E_{1012} \\
 E_{1020} & E_{1021} & E_{1022} \\
\end{array}
\right) & \left(
\begin{array}{ccc}
 E_{1100} & E_{1101} & E_{1102} \\
 E_{1110} & E_{1111} & E_{1112} \\
 E_{1120} & E_{1121} & E_{1122} \\
\end{array}
\right) & \left(
\begin{array}{ccc}
 E_{1200} & E_{1201} & E_{1202} \\
 E_{1210} & E_{1211} & E_{1212} \\
 E_{1220} & E_{1221} & E_{1222} \\
\end{array}
\right) \\
 \left(
\begin{array}{ccc}
 E_{2000} & E_{2001} & E_{2002} \\
 E_{2010} & E_{2011} & E_{2012} \\
 E_{2020} & E_{2021} & E_{2022} \\
\end{array}
\right) & \left(
\begin{array}{ccc}
 E_{2100} & E_{2101} & E_{2102} \\
 E_{2110} & E_{2111} & E_{2112} \\
 E_{2120} & E_{2121} & E_{2122} \\
\end{array}
\right) & \left(
\begin{array}{ccc}
 E_{2200} & E_{2201} & E_{2202} \\
 E_{2210} & E_{2211} & E_{2212} \\
 E_{2220} & E_{2221} & E_{2222} \\
\end{array}
\right) \\
\end{array}
\right)
$$
$$
\text{2D:}\quad
E_{ijkl} =
\left(
\begin{array}{cc}
 \left(
\begin{array}{cc}
 E_{0000} & E_{0001} \\
 E_{0010} & E_{0011} \\
\end{array}
\right) & \left(
\begin{array}{cc}
 E_{0100} & E_{0101} \\
 E_{0110} & E_{0111} \\
\end{array}
\right) \\
 \left(
\begin{array}{cc}
 E_{1000} & E_{1001} \\
 E_{1010} & E_{1011} \\
\end{array}
\right) & \left(
\begin{array}{cc}
 E_{1100} & E_{1101} \\
 E_{1110} & E_{1111} \\
\end{array}
\right) \\
\end{array}
\right)
$$
After accounting for the symmetries of an elasticity tensor (\ref{eqn:symmetry}),
these tensors really only have 21 and 6 unique entries in 3D and 2D respectively:
$$
\text{3D:}\quad
E_{ijkl} =
\left(\begin{array}{ccc}
 \left(\begin{array}{ccc}
 E_{0000} & E_{0001} & E_{0002} \\
 E_{0001} & E_{0011} & E_{0012} \\
 E_{0002} & E_{0012} & E_{0022} \\ \end{array} \right) &
 \left( \begin{array}{ccc}
 E_{0001} & E_{0101} & E_{0201} \\
 E_{0101} & E_{1101} & E_{1201} \\
 E_{0201} & E_{1201} & E_{2201} \\ \end{array} \right) &
 \left( \begin{array}{ccc}
 E_{0002} & E_{0201} & E_{0202} \\
 E_{0201} & E_{1102} & E_{1202} \\
 E_{0202} & E_{1202} & E_{2202} \\ \end{array} \right) \\
 \left( \begin{array}{ccc}
 E_{0001} & E_{0101} & E_{0201} \\
 E_{0101} & E_{1101} & E_{1201} \\
 E_{0201} & E_{1201} & E_{2201} \\ \end{array} \right) &
 \left( \begin{array}{ccc}
 E_{0011} & E_{1101} & E_{1102} \\
 E_{1101} & E_{1111} & E_{1112} \\
 E_{1102} & E_{1112} & E_{1122} \\ \end{array} \right) &
 \left( \begin{array}{ccc}
 E_{0012} & E_{1201} & E_{1202} \\
 E_{1201} & E_{1112} & E_{1212} \\
 E_{1202} & E_{1212} & E_{2212} \\ \end{array} \right) \\
 \left( \begin{array}{ccc}
 E_{0002} & E_{0201} & E_{0202} \\
 E_{0201} & E_{1102} & E_{1202} \\
 E_{0202} & E_{1202} & E_{2202} \\ \end{array} \right) &
 \left( \begin{array}{ccc}
 E_{0012} & E_{1201} & E_{1202} \\
 E_{1201} & E_{1112} & E_{1212} \\
 E_{1202} & E_{1212} & E_{2212} \\ \end{array} \right) &
 \left( \begin{array}{ccc}
 E_{0022} & E_{2201} & E_{2202} \\
 E_{2201} & E_{1122} & E_{2212} \\
 E_{2202} & E_{2212} & E_{2222} \\ \end{array} \right) \\
\end{array}
\right)
$$
$$
\text{2D:}\quad
E_{ijkl} =
\left(
\begin{array}{cc}
 \left(
\begin{array}{cc}
 E_{0000} & E_{0001} \\
 E_{0001} & E_{0011} \\
\end{array}
\right) & \left(
\begin{array}{cc}
 E_{0001} & E_{0101} \\
 E_{0101} & E_{1101} \\
\end{array}
\right) \\
 \left(
\begin{array}{cc}
 E_{0001} & E_{0101} \\
 E_{0101} & E_{1101} \\
\end{array}
\right) & \left(
\begin{array}{cc}
 E_{0011} & E_{1101} \\
 E_{1101} & E_{1111} \\
\end{array}
\right) \\
\end{array}
\right)
$$
And flattening to a symmetric matrix as described gets a single copy of each of
these entries in the upper triangle:
$$
\text{3D:}\quad
D_{\overline{ij}\, \overline{kl}} = 
\left(
\begin{array}{cccccc}
 E_{0000} & E_{0011} & E_{0022} & E_{0012} & E_{0002} & E_{0001} \\
 E_{0011} & E_{1111} & E_{1122} & E_{1112} & E_{1102} & E_{1101} \\
 E_{0022} & E_{1122} & E_{2222} & E_{2212} & E_{2202} & E_{2201} \\
 E_{0012} & E_{1112} & E_{2212} & E_{1212} & E_{1202} & E_{1201} \\
 E_{0002} & E_{1102} & E_{2202} & E_{1202} & E_{0202} & E_{0201} \\
 E_{0001} & E_{1101} & E_{2201} & E_{1201} & E_{0201} & E_{0101} \\
\end{array}
\right)
$$
$$
\text{2D:}\quad
D_{\overline{ij}\, \overline{kl}} = 
\left(
\begin{array}{ccc}
 E_{0000} & E_{0011} & E_{0001} \\
 E_{0011} & E_{1111} & E_{1101} \\
 E_{0001} & E_{1101} & E_{0101} \\
\end{array}
\right).
$$
Here $D$ is used to denote the matrix holding the flattened elasticity tensor,
as is standard in engineering literature.

\section{Implementing Double Contraction}
\subsection{The Missing Factor of 2}
We have flattened with the goal of using dot products to implement double contraction
of two rank 2 tensors and matrix-vector multiplication to implement double
contraction of a rank 4 tensor and a rank 2 tensor. However we're not done yet:
the off diagonal entries aren't handled properly! Consider a double contraction
of a stress tensor with a strain tensor in 3D:
$$
\left( \begin{array}{ccc}
 \sigma_{xx} & \sigma_{xy} & \sigma_{xz} \\
 \sigma_{xy} & \sigma_{yy} & \sigma_{yz} \\
 \sigma_{xz} & \sigma_{yz} & \sigma_{zz} \\
\end{array} \right) :
\left( \begin{array}{ccc}
 \epsilon_{xx} & \epsilon_{xy} & \epsilon_{xz} \\
 \epsilon_{xy} & \epsilon_{yy} & \epsilon_{yz} \\
 \epsilon_{xz} & \epsilon_{yz} & \epsilon_{zz} \\
\end{array} \right) = 
\sigma_{xx} \epsilon_{xx}+\sigma_{yy} \epsilon_{yy}+\sigma_{zz} \epsilon_{zz}+ 2\left(\sigma_{yz} \epsilon_{yz}+ \sigma_{xz} \epsilon_{xz} + \sigma_{xy} \epsilon_{xy} \right)
$$
Notice the factor of 2 on the off-diagonals, which isn't present in the dot product of the two tensors in flattened form:
$$
\begin{pmatrix}
    \sigma_{xx} &
    \sigma_{yy} &
    \sigma_{zz} &
    \sigma_{yz} &
    \sigma_{xz} &
    \sigma_{xy}
\end{pmatrix}
\begin{pmatrix}
    \e_{xx} \\
    \e_{yy} \\
    \e_{zz} \\
    \e_{yz} \\
    \e_{xz} \\
    \e_{xy}
\end{pmatrix} = \sigma_{xx} \epsilon_{xx}+\sigma_{yy} \epsilon_{yy}+\sigma_{zz} \epsilon_{zz}+ \sigma_{yz} \epsilon_{yz}+ \sigma_{xz} \epsilon_{xz} + \sigma_{xy} \epsilon_{xy}
$$
A similar problem happens with a double contraction of a rank 4 tensor with a
rank 2 tensor. Consider entries of the symmetric stress tensor $\sigma = E :
\epsilon$:
$$
\begin{pmatrix}
\sigma_{xx} \\
\sigma_{yy} \\
\sigma_{zz} \\
\sigma_{yz} \\
\sigma_{xz} \\
\sigma_{xy}
\end{pmatrix} =
\begin{pmatrix}
 E_{0000} \epsilon_{xx}+E_{0011} \epsilon_{yy}+E_{0022} \epsilon_{zz} +2 (E_{0012} \epsilon_{yz} + E_{0002} \epsilon_{xz}+ E_{0001} \epsilon_{xy}) \\
 E_{0011} \epsilon_{xx}+E_{1111} \epsilon_{yy}+E_{1122} \epsilon_{zz} +2 (E_{1112} \epsilon_{yz} + E_{1102} \epsilon_{xz}+ E_{1101} \epsilon_{xy}) \\
 E_{0022} \epsilon_{xx}+E_{1122} \epsilon_{yy}+E_{2222} \epsilon_{zz} +2 (E_{2212} \epsilon_{yz} + E_{2202} \epsilon_{xz}+ E_{2201} \epsilon_{xy}) \\
 E_{0012} \epsilon_{xx}+E_{1112} \epsilon_{yy}+E_{2212} \epsilon_{zz} +2 (E_{1212} \epsilon_{yz} + E_{1202} \epsilon_{xz}+ E_{1201} \epsilon_{xy}) \\
 E_{0002} \epsilon_{xx}+E_{1102} \epsilon_{yy}+E_{2202} \epsilon_{zz} +2 (E_{1202} \epsilon_{yz} + E_{0202} \epsilon_{xz}+ E_{0201} \epsilon_{xy}) \\
 E_{0001} \epsilon_{xx}+E_{1101} \epsilon_{yy}+E_{2201} \epsilon_{zz} +2 (E_{1201} \epsilon_{yz} + E_{0201} \epsilon_{xz}+ E_{0101} \epsilon_{xy})
\end{pmatrix}.
$$
Again, our flattened version of this operation (applying matrix $D$ to the flattened strain vector) gets almost the correct answer, but it misses
the factor of 2 on the off-diagonal strain contributions:
$$
D {\bf e} \defeq D
\begin{pmatrix}
\epsilon_{xx} \\
\epsilon_{yy} \\
\epsilon_{zz} \\
\epsilon_{yz} \\
\epsilon_{xz} \\
\epsilon_{xy}
\end{pmatrix} =
\begin{pmatrix}
E_{0000} \epsilon_{xx} + E_{0011} \epsilon_{yy} + E_{0022} \epsilon_{zz} + E_{0012} \epsilon_{yz} + E_{0002} \epsilon_{xz} + E_{0001} \epsilon_{xy} \\
E_{0011} \epsilon_{xx} + E_{1111} \epsilon_{yy} + E_{1122} \epsilon_{zz} + E_{1112} \epsilon_{yz} + E_{1102} \epsilon_{xz} + E_{1101} \epsilon_{xy} \\
E_{0022} \epsilon_{xx} + E_{1122} \epsilon_{yy} + E_{2222} \epsilon_{zz} + E_{2212} \epsilon_{yz} + E_{2202} \epsilon_{xz} + E_{2201} \epsilon_{xy} \\
E_{0012} \epsilon_{xx} + E_{1112} \epsilon_{yy} + E_{2212} \epsilon_{zz} + E_{1212} \epsilon_{yz} + E_{1202} \epsilon_{xz} + E_{1201} \epsilon_{xy} \\
E_{0002} \epsilon_{xx} + E_{1102} \epsilon_{yy} + E_{2202} \epsilon_{zz} + E_{1202} \epsilon_{yz} + E_{0202} \epsilon_{xz} + E_{0201} \epsilon_{xy} \\
E_{0001} \epsilon_{xx} + E_{1101} \epsilon_{yy} + E_{2201} \epsilon_{zz} + E_{1201} \epsilon_{yz} + E_{0201} \epsilon_{xz} + E_{0101} \epsilon_{xy}
\end{pmatrix}
$$
Finally, by symmetry, an identical problem happens when implementing the double
contraction $\e_{ij} E_{ijkl}$ as ${\bf e}^T D$.

\subsection{Solving the Problem}
There are several ways to introduce the factors of two. We can double the right
half of the D matrix, but this breaks its symmetry. We can leave $D$ unchanged,
but work with ``engineering strain'' $(\e_{xx}, \e_{yy}, \e_{zz}, 2 \e_{yz}, 2
\e_{xz}, 2 \e_{xy})^T$, but this means strain and stress are flattened
differently, which could lead to mistakes. We could weight the off diagonal
terms of all symmetric rank 2 tensors by $\sqrt{2}$ (this is called Kelvin or
Mandel notation) but this also requires modifying the $D$ matrix--though this
time in a symmetric way.

Instead of any of these, we choose to work with true, unmodified strain/stress
quantities and the original, symmetric $D$, but to include a "shear doubling"
matrix that doubles the off-diagonal entries each time a double contraction
is taken. This means off diagonals of strain and stress are flattened the
same way, $D$ remains symmetric, and entries in $D$ give the corresponding elasticity
coefficients directly instead of some scaling of them. This shear doubling
matrix is given by
$$
\text{3D: }\quad S = \begin{pmatrix}
        1 & 0 & 0 & 0 & 0 & 0 \\
        0 & 1 & 0 & 0 & 0 & 0 \\
        0 & 0 & 1 & 0 & 0 & 0 \\
        0 & 0 & 0 & 2 & 0 & 0 \\
        0 & 0 & 0 & 0 & 2 & 0 \\
        0 & 0 & 0 & 0 & 0 & 2 \\
    \end{pmatrix}, \quad \quad
\text{2D: }\quad S = \begin{pmatrix}
        1 & 0 & 0 \\
        0 & 1 & 0 \\
        0 & 0 & 2 \\
    \end{pmatrix}
$$
\subsection{Examples}
Finally, we show examples of double contraction implemented by matrix vector
multiplies and inner products. We use $F$ to denote the flattening operation
mapping symmetric rank 2 tensors to vectors and symmetric rank 4 tensors to
symmetric matrices. Letting ${\bf e} = F(\e)$ and ${\bf s} = F(\sigma)$
(flattened stress and strain tensors), we show various identities of our
notation:
\begin{gather*}
{\bf s} = F(\sigma) = F(E : \e) = F(E_{ijkl} \e_{kl}) = D S {\bf e}, \\
\sigma : \e = \sigma_{ij} \e_{kl} = {\bf s}^T S {\bf e}, \\
\e : E : \e = \e_{ij} E_{ijkl} \e_{kl} = {\bf e}^T S D S {\bf e}, \\
F(A : B) = F(A_{ijpq} B_{pqkl}) = F(A) S F(B)
\end{gather*}

Note, however, that the last operation does not result in a tensor with major
symmetries unless $A_{ijpq} = B_{pqij}$. This means that the matrix representing
the flattened result is not necessarily symmetric.

\section{Quadruple Contraction}
We can also implement quadruple contraction in flattened form:
$$
A_{ijkl} B_{ijkl}
$$
This is useful in several settings, including computing the Frobenius norm of a
tensor. Each term of the contraction will be an entry of F(A) multiplied by the
corresponding entry of F(B)---we just need to determine how many times each
flattened entry appears. That is, we can compute the contraction by instead sum over
entries of the flattened tensor, weighting each term by how many components are
flattened down to that particular entry. This can be determined from the minor
symmetries alone; major symmetry just causes a symmetric flattened tensor as
opposed to flattening more components of the tensor to a single flattened entry.

Below are the counts of how many tensor components are represented by each
flattened entry:
$$
\text{3D: }\begin{pmatrix}
        1 & 1 & 1 & 2 & 2 & 2 \\
        1 & 1 & 1 & 2 & 2 & 2 \\
        1 & 1 & 1 & 2 & 2 & 2 \\
        2 & 2 & 2 & 4 & 4 & 4 \\
        2 & 2 & 2 & 4 & 4 & 4 \\
        2 & 2 & 2 & 4 & 4 & 4 \\
    \end{pmatrix}, \quad \quad
\text{2D: }\begin{pmatrix}
        1 & 1 & 2 \\
        1 & 1 & 2 \\
        2 & 2 & 4 \\
    \end{pmatrix}
$$
This per-entry weighting can be applied to a flattened tensor by ``shear
doubling'' both the rows and columns,
$$
S F(A) S,
$$
allowing us to implement the quadruple contraction as a double contraction of
the scaled flattened tensors:
$$
A_{ijkl} B_{ijkl} = [S F(A) S] : F(B) = [F(A)S] : [S F(B)] = [SF(A)] : [F(B)S] =
S^{\frac{1}{2}} F(A) S^{\frac{1}{2}} : S^{\frac{1}{2}} F(B) S^{\frac{1}{2}} : 
$$
In linear algebra terms, this is:
$$
A_{ijkl} B_{ijkl} = \tr\left(S F(A)^T S F(B)\right) = \tr\left(F(A)^T S F(B) S\right)
= \tr\left(S^{\frac{1}{2}} F(A)^T S^{\frac{1}{2}} S^{\frac{1}{2}} F(B) S^{\frac{1}{2}} \right)
$$
\subsection{Examples}
As an example, we compute the squared Frobenius norm of an orthotropic elasticity
tensor:
$$
\text{3D: } \quad F(E) = \begin{pmatrix}
    E_{0000} & E_{0011} & E_{0022} & 0 & 0 & 0 \\
    E_{0011} & E_{1111} & E_{1122} & 0 & 0 & 0 \\
    E_{0022} & E_{1122} & E_{2222} & 0 & 0 & 0 \\
           0 &        0 &        0 & E_{1212} & 0 & 0 \\
           0 &        0 &        0 & 0 & E_{0202} & 0 \\
           0 &        0 &        0 & 0 & 0 & E_{0101} \\
\end{pmatrix}, \quad \quad
\text{2D: } \quad F(E) = \begin{pmatrix}
    E_{0000} & E_{0011} & 0 \\
    E_{0011} & E_{1111} & 0 \\
           0 &        0 & E_{0101} \\
\end{pmatrix}
$$
In 3D the Frobenius norm is
\begin{align*}
\norm{E}^2_F = &E_{ijkl} : E_{ijkl} = \norm{S^{1/2} F(E) S^{1/2}}^2_F = \\
&E_{0000}^2 + E_{1111}^2 + E_{2222}^2 + 2 E_{0011}^2 + 2 E_{0022}^2 + 2 E_{1122}^2
+ 4 E_{1212}^2 + 4 E_{0202}^2 + 4 E_{0101}^2,
\end{align*}
and in 2D:
\begin{align*}
\norm{E}^2_F = E_{ijkl} : E_{ijkl} = \norm{S^{1/2} F(E) S^{1/2}}^2_F = E_{0000}^2 + E_{1111}^2 + 2 E_{0011}^2 + 4 E_{0101}^2.
\end{align*}
Notice that the only difference between $\norm{E}_F^2$ and $\norm{F(E)}_F^2$
for orthotropic tensors is the weight of 4 on the ``shear'' terms.

\section{Rank 4 Identities and Inverses}
Defining the rank 4 identity and inverse operators in a way that is consistent
with our flattening and contraction conventions requires a bit of care. While
linear mappings on symmetric rank 2 tensors, $Sym^{n\times n} \to Sym^{n\times
n}$, do not have a unique representation as a rank 4 tensor, they do have a
unique representation with the symmetries required for flattening. We provide
this unique definition below.

\subsection{Symmetric Rank 4 Identity}
We flatten the identity tensor defined by its action  on $X \in Sym^{n \times
n}$:

\begin{equation}
\label{eqn:identityDef}
I_{ijkl} X_{kl} =  X_{ij}.
\end{equation}

Obviously the tensor $\delta_{ik} \delta_{jl}$ implements this map, but it
doesn't have the symmetries of an elasticity tensor and cannot be flattened.
However, since (\ref{eqn:identityDef}) does not
uniquely define I, we are free to choose the unique rank 4 tensor that has major and
minor symmetries:
$I_{ijkl} = \frac{1}{2}\left(\delta_{ik} \delta_{jl} + \delta_{il} \delta_{jk} \right)$.
This choice can be derived/motivated as follows.

We wish to find the symmetric tensor $I$ such that for all $X \in Sym^{nxn}$:
\begin{gather*}
F(I : X) = F(X) \quad \Longrightarrow \quad
F(I) S F(X) = F(X) \quad \Longrightarrow \quad
\boxed{
F(I) = S^{-1},}
\end{gather*}
which is indeed what we get when flattening the symmetrized identity,
$I_{ijkl} = \frac{1}{2}\left(\delta_{ik} \delta_{jl} + \delta_{il} \delta_{jk} \right)$. As
a sanity check, we can verify that this acts as an identity for double contraction of
rank 4 tensors as well:
$$
F(I : A) = F(I) S F(A) = S^{-1} S F(A) = F(A). \quad \checkmark
$$

\subsection{Symmetric Rank 4 Inverse}
Now we define the inverse, $A^{-1}$, of a symmetric rank 4 tensor $A$ by
\begin{equation}
\label{eqn:identityDef}
A_{ijpq} A^{-1}_{pqkl} = I_{ijkl},
\end{equation}
but again, due to symmetries, the definition is not unique. We derive the
unique inverse that has major and minor symmetries:
$$
F(A:A^{-1}) = F(I) \quad \Longrightarrow \quad  F(A) S F(A^{-1}) = S^{-1} \quad
\Longrightarrow \quad \boxed{F(A^{-1}) = S^{-1} F(A)^{-1} S^{-1},}
$$
and recover $A^{-1}$ itself by unflattening.

\subsection{Examples}
\subsubsection{Inverse of the Identity}
We do another simple sanity check of our notation, that the identity is its own
inverse:
$$
F(I^{-1}) = S^{-1} F(I)^{-1} S^{-1} = S^{-1} S S^{-1} = S^{-1} = F(I) \quad \checkmark
$$

\subsubsection{Orthotropic Material}
An orthotropic material parametrized by Young's moduli $Y$, Poisson ratios $\nu$
and shear moduli $\mu$ has a compliance tensor, $E^{-1}$ with a particularly simple
flattened form. The mapping from stress to strain tensors can be expressed as:
$$
\text{3D: } F(\epsilon) =
\begin{pmatrix}
    \frac{1}{Y_x} & -\frac{\nu_{yx}}{Y_y} & -\frac{\nu_{zx}}{Y_z} & 0 & 0 & 0 \\
    -\frac{\nu_{yx}}{Y_y} & \frac{1}{Y_y} & -\frac{\nu_{zy}}{Y_z} & 0 & 0 & 0 \\
    -\frac{\nu_{zx}}{Y_z} & -\frac{\nu_{zy}}{Y_z} & \frac{1}{Y_z} & 0 & 0 & 0 \\
    0 & 0 & 0 & \frac{1}{2 \mu_{yz}} & 0 & 0 \\
    0 & 0 & 0 & 0 & \frac{1}{2 \mu_{zx}} & 0 \\
    0 & 0 & 0 & 0 & 0 & \frac{1}{2 \mu_{xy}}
\end{pmatrix} F(\sigma), \quad
\text{2D: } F(\epsilon) = 
\begin{pmatrix}
    \frac{1}{Y_x} & -\frac{\nu_{yx}}{Y_y} & 0 \\
    -\frac{\nu_{yx}}{Y_y} & \frac{1}{Y_y}  & 0 \\
    0 & 0 & \frac{1}{2 \mu_{xy}}
\end{pmatrix} F(\sigma)
$$
Since our flattened double contraction formula inserts a shear doubler $S$ between its
arguments, we need to apply $S^{-1}$ to the matrix above to arrive at $F(E^{-1})$
(so that the double contraction $E^{-1} : \sigma$ is this mapping):
$$
\text{3D: } F(E^{-1}) =
\begin{pmatrix}
    \frac{1}{Y_x} & -\frac{\nu_{yx}}{Y_y} & -\frac{\nu_{zx}}{Y_z} & 0 & 0 & 0 \\
    -\frac{\nu_{yx}}{Y_y} & \frac{1}{Y_y} & -\frac{\nu_{zy}}{Y_z} & 0 & 0 & 0 \\
    -\frac{\nu_{zx}}{Y_z} & -\frac{\nu_{zy}}{Y_z} & \frac{1}{Y_z} & 0 & 0 & 0 \\
    0 & 0 & 0 & \frac{1}{4 \mu_{yz}} & 0 & 0 \\
    0 & 0 & 0 & 0 & \frac{1}{4 \mu_{zx}} & 0 \\
    0 & 0 & 0 & 0 & 0 & \frac{1}{4 \mu_{xy}}
\end{pmatrix}, \quad
\text{2D: } F(E^{-1}) = 
\begin{pmatrix}
    \frac{1}{Y_x} & -\frac{\nu_{yx}}{Y_y} & 0 \\
    -\frac{\nu_{yx}}{Y_y} & \frac{1}{Y_y}  & 0 \\
    0 & 0 & \frac{1}{4 \mu_{xy}}
\end{pmatrix}.
$$
Then we can apply our inverse formula, 
$
F(E) = S^{-1} F\left(E^{-1}\right)^{-1} S^{-1}$ to get:
\begin{gather*}
\text{3D: } F(E) =
\begin{pmatrix}
 \frac{Y_{x} Y_{y} \left(Y_{y} v_{zy}^2-Y_{z}\right)}{\text{denominator}} & -\frac{Y_{x} Y_{y} (Y_{z} v_{yx}+Y_{y} v_{zx} v_{zy})}{\text{denominator}} & -\frac{Y_{x} Y_{y} Y_{z} (v_{zx}+v_{yx} v_{zy})}{\text{denominator}} & 0 & 0 & 0 \\
 -\frac{Y_{x} Y_{y} (Y_{z} v_{yx}+Y_{y} v_{zx} v_{zy})}{\text{denominator}} & \frac{Y_{y}^2 \left(Y_{x} v_{zx}^2-Y_{z}\right)}{\text{denominator}} & -\frac{Y_{y} Y_{z} (Y_{x} v_{yx} v_{zx}+Y_{y} v_{zy})}{\text{denominator}} & 0 & 0 & 0 \\
 -\frac{Y_{x} Y_{y} Y_{z} (v_{zx}+v_{yx} v_{zy})}{\text{denominator}} & -\frac{Y_{y} Y_{z} (Y_{x} v_{yx} v_{zx}+Y_{y} v_{zy})}{\text{denominator}} & \frac{Y_{z}^2 \left(Y_{x} v_{yx}^2-Y_{y}\right)}{\text{denominator}} & 0 & 0 & 0 \\
 0 & 0 & 0 & \mu_{yz} & 0 & 0 \\
 0 & 0 & 0 & 0 & \mu_{zx} & 0 \\
 0 & 0 & 0 & 0 & 0 & \mu_{xy} \\
\end{pmatrix}, \\
\text{2D: } F(E) =
\begin{pmatrix}
    \frac{Y_{x} Y_{y}}{Y_{y}-Y_{x} \nu_{yx}^2} &
    \frac{Y_{x} Y_{y} \nu_{yx}}{Y_{y}-Y_{x} \nu_{yx}^2}
    & 0 \\
     \frac{Y_{x} Y_{y} \nu_{yx}}{Y_{y}-Y_{x} \nu_{yx}^2}
     & \frac{Y_{y}^2}{Y_{y}-Y_{x} \nu_{yx}^2} & 0 \\
      0 & 0 & \mu_{xy} \\
 \end{pmatrix}
\end{gather*}
with 3D denominator $Y_x Y_z \nu_{yx}^2 + Y_y^2 \nu_{zy}^2 + Y_y (-Y_z + Y_x
\nu_{zx} (\nu_{zx} + 2 \nu_{yx} \nu_{zy}))$.

\end{document}
